\documentclass{article}
\usepackage{amsmath,amssymb,amsthm,mathtools}
\newtheorem{theorem}{Theorem}
\newtheorem{lemma}{Lemma}
\newtheorem{proposition}{Proposition}
\newtheorem{corollary}{Corollary}
\newtheorem{definition}{Definition}
\newtheorem{example}{Example}
\title{ShadowBench geometry/L3/geo\_gen\_L3\_001}
\begin{document}
\maketitle
% Problem id: geometry/L3/geo_gen_L3_001
% Companion instructions: docs/instructions.md
% Candidate skeletons: docs/skeletons/
% Lean target: ShadowBench/Source/Main.lean

\begin{theorem}[morleys_trisector_theorem]\label{thm:morley_trisector}
In plane geometry, Morley's trisector theorem states that in any triangle, the three points of intersection of the adjacent angle trisectors form an equilateral triangle, called the \textbf{first Morley triangle}.

Formally, let $X, Y, Z$ be the intersections of the adjacent internal angle trisectors of a triangle $ABC$. Then the triangle $XYZ$ is equilateral.
\end{theorem}
\end{document}
